\documentclass[../../main]{subfiles}

\begin{document}

\section{Cardinalidade de Espaços de Funções}
\label{sec:matematica:funcoes:cardinalidade}

\subsection{Número de Funções}

\begin{prop}
  \label{prop:matematica:funcoes:cardinalidade:numero-funcoes}

  Sejam \(A\) e \(B\) conjuntos finitos.

  Então o número de funções

  \[
    f:A\to B
  \]

  é dado por

  \[
    |B|^{|A|}.
  \]

\end{prop}

\begin{proof}

  Para cada elemento

  \[
    a\in A
  \]

  existem

  \[
    |B|
  \]

  escolhas possíveis para o valor

  \[
    f(a).
  \]

  Como as escolhas são independentes para cada elemento de \(A\),
  pelo princípio multiplicativo,

  \[
    |B^A|
    =
    |B|^{|A|}.
  \]

\end{proof}

\begin{ex}

  O número de funções

  \[
    \{1,2,3\}
    \to
    \{0,1\}
  \]

  é

  \[
    2^3
    =
    8.
  \]

\end{ex}

\subsection{Número de Funções Injetivas}

\begin{prop}
  \label{prop:matematica:funcoes:cardinalidade:injetivas}

  Sejam \(A\) e \(B\) conjuntos finitos com

  \[
    |A|=n
  \]

  e

  \[
    |B|=m.
  \]

  Se

  \[
    n\le m,
  \]

  então o número de funções injetivas

  \[
    A\to B
  \]

  é

  \[
    m(m-1)\cdots(m-n+1).
  \]

\end{prop}

\begin{proof}

  O primeiro elemento de \(A\) possui \(m\) possíveis imagens.

  O segundo possui \(m-1\).

  O terceiro possui \(m-2\).

  Prosseguindo dessa forma,

  \[
    m(m-1)\cdots(m-n+1)
  \]

  funções injetivas são obtidas.

\end{proof}

\begin{cor}

  O número de injeções

  \[
    A\to B
  \]

  é

  \[
    \frac{m!}{(m-n)!}.
  \]

\end{cor}

\subsection{Número de Funções Sobrejetivas}

\begin{prop}
  \label{prop:matematica:funcoes:cardinalidade:sobrejetivas}

  Sejam \(A\) e \(B\) conjuntos finitos com

  \[
    |A|=n,
    \qquad
    |B|=m.
  \]

  O número de sobrejeções

  \[
    A\to B
  \]

  é

  \[
    \sum_{k=0}^{m}
    (-1)^k
    \binom{m}{k}
    (m-k)^n.
  \]

\end{prop}

\begin{proof}

  Aplicação do princípio da inclusão-exclusão.

\end{proof}

\begin{obs}

  A demonstração completa requer técnicas de análise combinatória e
  será apresentada posteriormente.

\end{obs}

\subsection{Número de Bijeções}

\begin{prop}
  \label{prop:matematica:funcoes:cardinalidade:bijecoes}

  Se

  \[
    |A|
    =
    |B|
    =
    n,
  \]

  então o número de bijeções

  \[
    A\to B
  \]

  é

  \[
    n!.
  \]

\end{prop}

\begin{proof}

  Toda bijeção é uma função injetiva.

  Aplicando o resultado anterior com

  \[
    m=n,
  \]

  obtemos

  \[
    n(n-1)\cdots 2\cdot1
    =
    n!.
  \]

\end{proof}

\begin{cor}

  O número de permutações de um conjunto finito com \(n\) elementos é

  \[
    n!.
  \]

\end{cor}

\subsection{Funções Características}

\begin{prop}
  \label{prop:matematica:funcoes:cardinalidade:caracteristicas}

  Seja \(X\) um conjunto finito.

  Então

  \[
    |\mathcal P(X)|
    =
    |\{0,1\}^X|
    =
    2^{|X|}.
  \]

\end{prop}

\begin{proof}

  Pela bijeção

  \[
    A
    \mapsto
    \chi_A,
  \]

  existe uma correspondência biunívoca entre

  \[
    \mathcal P(X)
  \]

  e

  \[
    \{0,1\}^X.
  \]

  Como

  \[
    |\{0,1\}|
    =
    2,
  \]

  segue que

  \[
    |\{0,1\}^X|
    =
    2^{|X|}.
  \]

\end{proof}

\subsection{Cardinalidade de Famílias Indexadas}

\begin{prop}
  \label{prop:matematica:funcoes:cardinalidade:familias}

  Seja

  \[
    I
  \]

  um conjunto finito e, para cada

  \[
    i\in I,
  \]

  seja dado um conjunto finito \(A_i\).

  Então

  \[
    \left|
    \prod_{i\in I}
    A_i
    \right|
    =
    \prod_{i\in I}
    |A_i|.
  \]

\end{prop}

\begin{proof}

  Uma família

  \[
    (a_i)_{i\in I}
  \]

  corresponde à escolha independente de um elemento em cada conjunto
  \(A_i\).

  Aplicando o princípio multiplicativo obtém-se o resultado.

\end{proof}

\subsection{Princípio das Gavetas}

\begin{thm}[Princípio das Gavetas]
  \label{thm:matematica:funcoes:cardinalidade:gavetas}

  Se

  \[
    |A|
    >
    |B|,
  \]

  então não existe função injetiva

  \[
    A\to B.
  \]

\end{thm}

\begin{proof}

  Caso contrário existiria uma injeção de um conjunto maior em um
  conjunto menor, contradizendo a contagem obtida anteriormente.

\end{proof}

\begin{cor}

  Se

  \[
    f:A\to B
  \]

  e

  \[
    |A|>|B|,
  \]

  então existem

  \[
    x\neq y
  \]

  tais que

  \[
    f(x)=f(y).
  \]

\end{cor}

\end{document}
